
\documentclass[tikz,border=8pt]{standalone}
\usepackage[UTF8,fontset=fandol]{ctex}
\usepackage{amsmath,amssymb}
\usetikzlibrary{arrows.meta,calc,positioning}
\definecolor{ink}{HTML}{172033}
\definecolor{muted}{HTML}{637083}
\definecolor{blue}{HTML}{4F86C6}
\definecolor{green}{HTML}{61A66A}
\definecolor{red}{HTML}{D55E5E}
\definecolor{amber}{HTML}{D59B3D}
\definecolor{paperblue}{HTML}{F2F7FD}
\definecolor{papergreen}{HTML}{F3F9F1}
\definecolor{paperamber}{HTML}{FFF8E8}
\tikzset{
  panel/.style={anchor=north west,minimum width=9.15cm,minimum height=6.05cm,
    text width=8.45cm,align=left,inner sep=10pt,draw=black!12,rounded corners=5pt,fill=white},
  title/.style={font=\bfseries\large,text=ink,anchor=west},
  body/.style={font=\small,text=muted,align=left,text width=7.9cm},
  formula/.style={draw=blue!35,fill=paperblue,rounded corners=4pt,inner sep=8pt,
    text=ink,align=center,text width=7.7cm},
  insight/.style={draw=green!40,fill=papergreen,rounded corners=4pt,inner sep=7pt,
    text=ink,align=left,text width=7.7cm,font=\small},
  arrow/.style={-{Latex[length=2.3mm]},line width=1pt,draw=blue!70},
  chip/.style={rounded corners=3pt,draw=black!14,fill=black!2,inner xsep=7pt,
    inner ysep=5pt,font=\small,text=ink},
}
\newcommand{\Panel}[5]{%
  \node[panel] (#1) at (#2,#3) {};
  \node[title] at ($(#1.north west)+(0.35,-0.42)$) {#4};
  \node[body,anchor=north west] at ($(#1.north west)+(0.35,-1.05)$) {#5};
}
\newcommand{\Summary}[1]{%
  \node[draw=amber!35,fill=paperamber,rounded corners=6pt,minimum width=28.2cm,
    minimum height=1.55cm,text width=27.2cm,align=center,font=\small,text=ink]
    at (14.125,-13.35) {\textbf{一图总结}\quad #1};
}
\newcommand{\Identity}{%
  \node[font=\scriptsize,text=black!38,anchor=east] at (28.15,-14.35)
    {cs229-storyboard@五道口纳什};
}
\newcommand{\Formula}[1]{%
  \par\vspace{5pt}\begingroup\setlength{\fboxsep}{8pt}%
  \colorbox{paperblue}{\parbox{7.45cm}{\centering\color{ink}#1}}%
  \endgroup\par\vspace{6pt}}
\newcommand{\Insight}[1]{%
  \par\vspace{5pt}\begingroup\setlength{\fboxsep}{7pt}%
  \colorbox{papergreen}{\parbox{7.55cm}{\color{ink}#1}}%
  \endgroup\par\vspace{6pt}}
\newcommand{\Chip}[1]{%
  \begingroup\setlength{\fboxsep}{5pt}\colorbox{black!4}{\strut\color{ink}#1}\endgroup}

\begin{document}
\begin{tikzpicture}[x=1cm,y=1cm]
\fill[white] (-0.35,0.35) rectangle (28.6,-14.55);
\Panel{p1}{0}{0}{1. 把预测写成矩阵}{设计矩阵把所有样本排成行：\Formula{$\hat y=X\theta$\\$X:m\times n,\ \theta:n\times1$}一次矩阵乘法得到全部预测。}
\Panel{p2}{9.55}{0}{2. 残差变成目标函数}{误差向量为 $r=X\theta-y$。\Formula{$J(\theta)=\tfrac12\lVert X\theta-y\rVert_2^2$}平方损失让大误差付出更高代价。}
\Panel{p3}{19.1}{0}{3. 求梯度并更新}{链式法则把残差拉回参数空间：\Formula{$\nabla_\theta J=X^\top(X\theta-y)$\\$\theta\leftarrow\theta-\alpha\nabla J$}}
\Panel{p4}{0}{-6.45}{4. 也可以直接求解}{令梯度为零：\Formula{$X^\top X\theta=X^\top y$\\$\theta^*=(X^\top X)^{-1}X^\top y$}实际计算优先用线性求解器或伪逆。}
\Panel{p5}{9.55}{-6.45}{5. 几何上是投影}{最优预测 $X\theta^*$ 是 $y$ 在 $X$ 列空间上的投影。\Insight{$X^\top(y-X\theta^*)=0$：残差与每个特征方向正交。}}
\Panel{p6}{19.1}{-6.45}{6. 两种方法怎么选}{\Chip{小而稠密：正规方程}\quad\Chip{大而稀疏：梯度法}无论选哪种，都先检查维度、缩放和验证集误差。}
\Summary{数据 $X,y$ $\rightarrow$ 预测 $X\theta$ $\rightarrow$ 残差 $r$ $\rightarrow$ 梯度或正规方程 $\rightarrow$ 最小二乘解。}
\Identity
\end{tikzpicture}
\end{document}
